\lecture{4}{На отл(8-9)}

\section{На отл(8-9)}

\subsection{Шаблоны для компилятора}
\subsubsection{Инстанцирование шаблонов}
\subsubsection{Правила вывода шаблонных типов}
\subsubsection{Вычисления на шаблонах: Фибоначчи, проверка простоты}
\subsubsection{Fold-expressions: 4 вида, синтаксис, раскрытие}
\subsubsection{Пример: \texttt{print} от переменного числа аргументов}
\subsubsection{Метафункции с fold: все типы одинаковы и т.п.}

\subsection{Виртуальное наследование}
\subsubsection{Расположение объектов в памяти}
\subsubsection{Устройство виртуального наследования (точка зрения компилятора)}
\subsubsection{Устройство виртуальных функций}
\subsubsection{Случай с множественным наследованием}

\subsection{Аллокаторы: нюансы}
\subsubsection{\texttt{propagate\_on\_container\_copy\_assignment}}
\subsubsection{\texttt{propagate\_on\_container\_move\_assignment}}
\subsubsection{\texttt{select\_on\_container\_copy\_construction}}
\subsubsection{Rebind аллокаторов}
\subsubsection{Реализация: конструктор копирования, копирующее и перемещающее присваивание для \texttt{list}}

\subsection{\texttt{std::forward}}
\subsubsection{Реализация и детальный разбор всех случаев}

\subsection{Вывод типов}
\subsubsection{\texttt{decltype} и отличия от \texttt{auto}}
\subsubsection{\texttt{decltype(auto)}}
\subsubsection{\texttt{auto} в аргументах функции и шаблонных аргументах}
\subsubsection{Structured bindings: как заставить работать со своим классом}
\subsubsection{\texttt{std::tuple}: пользование и идея реализации}

\subsection{Умные указатели и наследование}
\subsubsection{Проблема указателя на наследника}
\subsubsection{CRTP}
\subsubsection{\texttt{enable\_shared\_from\_this}: реализация}
\subsubsection{Проблема кастомных \texttt{deleter} и \texttt{allocator}}
\subsubsection{Type-erasure, \texttt{std::any}: идея реализации}
\subsubsection{Решение проблем с кастомными \texttt{deleter}/\texttt{allocator}}

\subsection{Type-traits: конструктор по умолчанию}
\subsubsection{\texttt{declval}}
\subsubsection{\texttt{move\_if\_noexcept}: идея реализации}
\subsubsection{\texttt{std::is\_base\_of}: реализация}

\subsection{Union и вариантные типы}
\subsubsection{Union}
\subsubsection{SSO в \texttt{std::string} с помощью \texttt{union}}
\subsubsection{\texttt{std::variant}: идея реализации}

\subsection{Ranges}
\subsubsection{Причины появления и примеры использования}
\subsubsection{Концепт range}
\subsubsection{\texttt{ranges::sort}}
\subsubsection{\texttt{borrowed\_iterator}}
\subsubsection{views, sentinel, \texttt{subrange} и \texttt{common\_view}}
\subsubsection{Комбинация отображений}
\subsubsection{Устройство \texttt{transform\_view}}

\subsection{\texttt{constexpr}}
\subsubsection{\texttt{constexpr} переменные}
\subsubsection{\texttt{constexpr} функции}
\subsubsection{\texttt{constexpr} в классах}
